\section{Additive characters and Frobenius detection}\label{sec:characters}

We will compare homomorphisms of Galois groups by evaluating additive
$\ell$-adic characters at Frobenius elements.  The main fact needed for
this comparison is that infinitely many primes of absolute residue
degree one detect a character.  We first prove the logarithm-rank
statement behind this fact, then establish the versions involving
averaging and corestriction.  The section ends with a representation
contained in the character space that will supply enough characters for
the comparison.

\subsection{A logarithm-rank lemma}

For a finite extension $K/\Q_\ell$, normalize its absolute value by
$|\ell|_\ell=\ell^{-1}$.  The $\ell$-adic logarithm on
$\mathcal O_K^\times$ is the continuous homomorphism that agrees with
the usual power series near $1$; its kernel is the finite group of
roots of unity in $K$.  Logarithms of points in a torus are taken
coordinatewise.

For a number field $M$, class field theory expresses an additive
character of $G_M$ as a linear form on $M\otimes_\Q\Q_\ell$.
A zero Frobenius value at a prime $v\nmid\ell$ forces this form to
annihilate $\log b$, where $(b)=v^h$ for some $h>0$.
For primes of residue degree one chosen over distinct rational primes
unramified in a Galois closure, the conjugate generators have disjoint
ideal supports; the following lemma turns their multiplicative
independence into full logarithmic rank.

This special rank estimate admits an
interpolation-determinant proof.  It is a special case of the
$\ell$-adic Waldschmidt--Masser rank theorem
\cite{waldschmidt1981,masser1981}.  We retain a proof using the
translate estimates of that theory and Laurent's interpolation
determinants \cite{laurent1991}; see also
\cite[Theorems~5.1 and~5.8]{dasgupta2023}.

\begin{lemma}[Logarithm rank]\label{lem:log-rank}
Let $n,m\geq1$, let $E$ be a number field embedded in a finite
extension $K/\Q_\ell$, and let
\[
 u_i=(u_{i1},\ldots,u_{in})\in(E^\times)^n,
 \qquad 1\leq i\leq m,
\]
have algebraic-integer coordinates that are units in $K$.  Suppose that
for every $a=(a_1,\ldots,a_m)\in\Z^m\setminus\{0\}$, the coordinates
of $u^a:=\prod_{i=1}^m u_i^{a_i}$ are multiplicatively independent
modulo roots of unity.  If $m>n(n-1)$, then
\[
 \operatorname{span}_K\{\log u_1,\ldots,\log u_m\}=K^n.
\]
\end{lemma}

\begin{proof}
The multiplicative hypothesis implies that the cyclic subgroup
generated by every $u^a$ with $a\ne0$ is Zariski dense in
$\mathbb G_m^n$.  Indeed, distinct Laurent monomials have distinct
values at $u^a$.  If a Laurent polynomial vanished at all nonnegative
powers of $u^a$, a Vandermonde matrix applied to its finitely many
monomials would force all its coefficients to vanish.

Suppose that the logarithms span a subspace of dimension $r<n$.
The case $r=0$ contradicts the multiplicative hypothesis because the
kernel of the logarithm on local units consists of roots of unity.
Thus $1\leq r<n$, and in particular $n\geq2$.  We will obtain an
algebraic determinant whose small $\ell$-adic absolute value is
incompatible with its archimedean size.

First make all choices needed for the analytic estimate.  Replacing
all $u_i$ by the same positive integer power preserves the hypotheses
and the logarithmic rank.  Such a power puts their coordinates in a
neighborhood where exponential and logarithm are inverse.  The
$\mathcal O_K$-module generated by their logarithms is free of rank
$r$; choose a basis $b_1,\ldots,b_r\in K^n$.  After a further common
$\ell$-power, these basis vectors have sufficiently small coordinates
that
\[
 \rho:=\max_{j,t}|b_{t,j}|_\ell
 \quad\hbox{satisfies}\quad
 \eta:=\rho\ell^{1/(\ell-1)}<1.
\]
Here the basis is multiplied by the same power as the logarithms.
Consequently there are $z_i\in\mathcal O_K^r$ such that
\[
 \log u_i=\sum_{t=1}^r z_{i,t}b_t.
\]
From now on these powered points, the number field $E$, the local
field $K$, the basis, and $\eta$ are fixed.  In particular, none
depends on the degree parameter introduced below.

For an integer vector $\alpha=(\alpha_1,\ldots,\alpha_n)$, the
monomial $X^\alpha$ on these points is represented by
\[
 f_\alpha(z)=
 \exp\!\left(\sum_{j=1}^n\alpha_j
                      \sum_{t=1}^r b_{t,j}z_t\right).
\]
It is analytic on a polydisc of radius greater than $1$.  If
$f_\alpha(z)=\sum_{\nu\in\Z_{\geq0}^r}c_{\alpha,\nu}z^\nu$,
the exponential series and
$|k!|_\ell^{-1}\leq\ell^{k/(\ell-1)}$ give the uniform bound
\begin{equation}\label{eq:log-taylor-bound}
 |c_{\alpha,\nu}|_\ell\leq\eta^{|\nu|},
 \qquad |\nu|=\nu_1+\cdots+\nu_r.
\end{equation}
The bound is independent of $\alpha$, since every integer has
$\ell$-adic absolute value at most $1$.  Moreover, for every
$a\in\Z_{\geq0}^m$, the point
$z(a):=\sum_i a_i z_i$ lies in $\mathcal O_K^r$, and
$f_\alpha(z(a))=(u^a)^\alpha$.

We next show that sufficiently many of these evaluations have full
algebraic rank.  For an integer $D\geq1$, put $A=(D+1)^n$ and choose
an integer $S\geq1$ such that
\begin{equation}\label{eq:log-grid-size}
 (S+1)^m>(nD)^n,
 \qquad S=O(D^{n/m}).
\end{equation}
Consider the matrix with rows indexed by
$\alpha\in\{0,\ldots,D\}^n$, columns indexed by
$a\in\{0,\ldots,nS\}^m$, and entries $(u^a)^\alpha$.
We claim that its rank is $A$.

Otherwise a nonzero polynomial $P$ of degree at most $D$ in each
variable would vanish at every point $u^a$ in this grid.  Over an
algebraic closure of $E$, define closed subsets of the torus by
\[
 Z_j=\{x\in\mathbb G_m^n:
          P(u^a x)=0\ \hbox{for all }a\in\{0,\ldots,jS\}^m\},
 \qquad 0\leq j\leq n.
\]
These sets are nested and nonempty, since $1\in Z_n$, whereas
$\dim Z_0\leq n-1$.  Hence some adjacent pair $Z_j,Z_{j+1}$ has
the same dimension $d$.  Let $V$ be a $d$-dimensional irreducible
component of $Z_{j+1}$.  For every
$b\in\{0,\ldots,S\}^m$, the translate $u^bV$ is contained in
$Z_j$ and has dimension $d$, so it is an irreducible component of
$Z_j$.  All these translates are distinct: an equality
$u^bV=u^{b'}V$ for $b\ne b'$ would make $V$ stable under every
power of $u^{b-b'}$.  The density proved above would then force
$V=\mathbb G_m^n$, contrary to $V\subseteq Z_0$.

For completeness, a common zero set of polynomials of total degree
at most $\delta\geq1$ in $\mathbb A^n$ has at most $\delta^n$
irreducible components.  One can obtain this form of the affine
B\'ezout bound by assigning a component $Y$ the weight
$\delta^{\dim Y}\deg Y$.  Start with $\mathbb A^n$, of weight
$\delta^n$, and impose the equations successively.  A component
contained in the next hypersurface retains its weight; a component
not contained in it is cut into components of dimension one less,
whose total degree is at most $\delta\deg Y$ by B\'ezout.  Thus
the total weight does not increase, and removing redundant components
can only decrease it.  The same component bound holds after
restriction to the torus.  Applying it to $Z_j$, whose equations
have total degree at most $nD$, contradicts
\eqref{eq:log-grid-size}.  This proves the rank claim.

Choose a nonzero $A\times A$ minor $\Delta_D$ of the evaluation
matrix.  It is an algebraic integer in the fixed field $E$.
Expand its rows using the convergent power series for $f_\alpha$.
A nonzero term must involve $A$ distinct monomials in the $r$
variables $z_1,\ldots,z_r$, since repeated monomials give repeated
rows in the corresponding evaluation determinant.  The sum of the
degrees of any $A$ distinct such monomials is at least
$c_rA^{1+1/r}$ for all sufficiently large $A$: the number of
monomials of degree at most $t$ is $\binom{t+r}{r}=O_r(t^r)$,
so at least half of the $A$ monomials have degree bounded below by
a positive multiple of $A^{1/r}$.  Their evaluations have absolute
value at most $1$.  The nonarchimedean triangle inequality and
\eqref{eq:log-taylor-bound}, first for truncated series and then by
convergence, therefore give a constant $c>0$ such that
\begin{equation}\label{eq:log-determinant-small}
 \log|\Delta_D|_\ell\leq-cA^{1+1/r}.
\end{equation}

At every archimedean embedding $\tau:E\hookrightarrow\C$, an
entry of this minor has absolute value at most $\exp(CDS)$, with
$C$ independent of $D$.  Expanding the determinant thus gives
\begin{equation}\label{eq:log-determinant-large}
 \log|\tau(\Delta_D)|\leq C'ADS+A\log A.
\end{equation}
Let $\lambda$ be the place of $E$ induced by $E\hookrightarrow K$.
Use absolute values at finite places extending the usual ones on
$\Q$, and sum over all embeddings $E\hookrightarrow\C$, so that
each complex place occurs twice.  Writing $p_v$ for the rational
prime below a finite place $v$, the product formula reads
\[
 0=\sum_{\tau:E\hookrightarrow\C}\log|\tau(\Delta_D)|
     +\sum_{v\text{ finite}}[E_v:\Q_{p_v}]
                            \log|\Delta_D|_v.
\]
Every term in the second sum is nonpositive because $\Delta_D$ is
an algebraic integer.  Retaining its $\lambda$-term and using
\eqref{eq:log-determinant-small}--\eqref{eq:log-determinant-large}
would imply
\[
 c[E_\lambda:\Q_\ell]A^{1+1/r}
 \leq [E:\Q](C'ADS+A\log A).
\]
This is impossible as $D\to\infty$: the left side has order
$D^{n+n/r}$, whereas the right side is
$O(D^{n+1+n/m}+D^n\log D)$, and
\[
 \frac nr\geq\frac n{n-1}>1+\frac nm
\]
by $m>n(n-1)$.  This contradiction proves the lemma.
\end{proof}

\subsection{Detection, averaging, and corestriction}

For a number field $M\subset\overline{\Q}$, write
\[
 \X{\ell}{M}:=\Hom_{\mathrm{cont}}(G_M,\Q_\ell),
 \qquad G_M=\Gal(\overline{\Q}/M),
\]
where $\Q_\ell$ is an additive group with trivial Galois action.
Local class field theory shows that these characters are unramified
at every finite prime $v\nmid\ell$: the local unit group has a
pro-$p$ subgroup of finite index for $p\ne\ell$, and has no
nonzero continuous homomorphism to $\Q_\ell$.  Thus
$x(\Fr_v)$ is defined for $x\in\X{\ell}{M}$ and $v\nmid\ell$.

If $M/\Q$ is Galois, $R=\Gal(M/\Q)$ acts by
$(g x)(h)=x(\widetilde g^{-1}h\widetilde g)$, where
$\widetilde g\in G_\Q$ lifts $g$.  The definition is independent
of the lift.  For a subgroup $J\subseteq R$, write
\[
 E_J=\frac1{|J|}\sum_{g\in J}g
\]
for the projector onto the $J$-invariants.  We use
$\res_{L/M}$ and $\cor_{L/M}$ for restriction and corestriction
on additive characters for a finite extension $L/M$.

\begin{proposition}[Frobenius detection]\label{prop:character-detection}
Let $M$ be a number field, let $\ell$ be a rational prime, and let
$x\in\X{\ell}{M}$.
\begin{enumerate}
\item[(i)] If $x(\Fr_v)=0$ at infinitely many primes $v\nmid\ell$
of absolute residue degree one, then $x=0$.
\item[(ii)] Suppose that $M/\Q$ is Galois, and fix
$\delta\in\Gal(M/\Q)$.  If $x(\Fr_v)=0$ at infinitely many
primes $v\nmid\ell$ lying over rational primes unramified in $M$
whose Frobenius element at $v$ in $\Gal(M/\Q)$ is $\delta$, then
$E_{\langle\delta\rangle}x=0$.
\item[(iii)] If $\cor_{M/\Q}x\ne0$, then
$x(\Fr_v)=0$ at only finitely many primes $v\nmid\ell$, with no
restriction on their residue degrees.
\end{enumerate}
\end{proposition}

\begin{proof}
We first recall the class-field-theoretic description of the character
space.  Let $V_M=M\otimes_\Q\Q_\ell$, and let $U_M\subseteq V_M$
be the $\Q_\ell$-linear span of the logarithms of the global units,
with logarithms taken at all places above $\ell$.  Then
\begin{equation}\label{eq:character-class-field}
 \X{\ell}{M}\simeq
 \{\lambda\in\Hom_{\Q_\ell}(V_M,\Q_\ell):
                         \lambda(U_M)=0\}.
\end{equation}
Indeed, pass to the inverse limit of the ray class sequences with
moduli supported above $\ell$.  Local logarithms identify the
rationalized local unit groups with $V_M$; the global units impose
exactly the subspace $U_M$.  The ordinary class group and the factors
at real places are finite, so contribute no additive
$\Q_\ell$-characters.  This description uses the actual logarithmic
image of the units and makes no assertion about its dimension.

For $v\nmid\ell$, choose $h_v>0$ and $b_v\in\mathcal O_M$ with
$(b_v)=v^{h_v}$.  Such a choice is possible because the ideal class
group is finite.  Principal reciprocity, with arithmetic Frobenius,
gives for the linear form corresponding to $x$ the identity
\begin{equation}\label{eq:character-principal-reciprocity}
 \lambda(\log b_v)=-h_vx(\Fr_v).
\end{equation}
Here $b_v$ is a unit at every place above $\ell$; all other local
unit contributions are killed by $x$.

For (i), set $n=[M:\Q]$ and choose $m>n(n-1)$ of the zero primes,
over distinct rational primes unramified in a Galois closure $E$ of
$M$.  Let $b_i$ be the elements just chosen for these primes, and
let
\[
 u_i=(\sigma(b_i))_{\sigma:M\hookrightarrow E}\in(E^\times)^n.
\]
All coordinates are algebraic integers and are units at every place
of $E$ above $\ell$.  We verify the multiplicative hypothesis of
Lemma~\ref{lem:log-rank} by valuations.

Put $G=\Gal(E/\Q)$ and $H=\Gal(E/M)$.  For a selected degree-one
prime $v_i$ and a prime $\mathfrak p_i$ of $E$ above it, the
decomposition group $D_i$ is contained in $H$.  In fact, the rational
prime is unramified in $E$, and the degree-one condition says that
$[D_i:D_i\cap H]=1$.  The support of $v_i\mathcal O_E$ is
indexed by $H/D_i$.  Its translates corresponding to the embeddings
$M\hookrightarrow E$, indexed by $G/H$, have supports indexed by
the disjoint sets $gH/D_i$.  Distinct indices $i$ involve distinct
rational primes.  It follows that the $mn$ numbers
$\sigma(b_i)$ have pairwise disjoint nonempty supports of their
principal ideals.  They are therefore multiplicatively independent
modulo roots of unity.  In particular, for any nonzero
$a\in\Z^m$, the $n$ coordinates of $u^a$ are multiplicatively
independent modulo roots of unity.

Embed $E$ in a finite extension $K/\Q_\ell$.  Under the splitting
\[
 (M\otimes_\Q\Q_\ell)\otimes_{\Q_\ell}K
   \simeq\prod_{\sigma:M\hookrightarrow E}K,
\]
the vector $\log b_i$ becomes $\log u_i$.
Lemma~\ref{lem:log-rank} says that these vectors span the right
side.  Equations~\eqref{eq:character-class-field} and
\eqref{eq:character-principal-reciprocity} therefore force
$\lambda=0$, and hence $x=0$.

For (ii), put $C=\langle\delta\rangle$, $F=M^C$, and
$y=E_Cx$.  Since $\delta$ fixes every indicated prime $v$, averaging
preserves its Frobenius value, so $y(\Fr_v)=0$.  The
restriction--corestriction identity gives
\[
 \res_{M/F}\cor_{M/F}y=\sum_{g\in C}g y=|C|y.
\]
Thus $y$ is the restriction of
$z=|C|^{-1}\cor_{M/F}y\in\X{\ell}{F}$.  At an indicated prime,
the decomposition group of $M/\Q$ is exactly $C$.  The prime
$u=v\cap F$ consequently has absolute residue degree one, and
\[
 0=y(\Fr_v)=[\kappa(v):\kappa(u)]z(\Fr_u).
\]
Infinitely many distinct $u$ occur because $M/F$ is finite.
Part~(i) gives $z=0$, and therefore $y=0$.

For (iii), suppose that infinitely many zero primes exist.  Let
$E/\Q$ be a finite Galois closure of $M$, put
$R=\Gal(E/\Q)$, and let $y=\res_{E/M}x$.  Discarding finitely
many primes, choose primes of $E$ above these zero primes whose
rational primes are unramified in $E$.  Restriction multiplies a
Frobenius value by the relative residue degree, so they are zero
primes for $y$.  An infinite subset has one fixed Frobenius element
$\delta\in R$.  By (ii),
$E_{\langle\delta\rangle}y=0$, and hence $E_Ry=0$.  It follows that
\[
 \res_{E/\Q}\cor_{E/\Q}y=|R|E_Ry=0.
\]
Restriction on additive characters across a finite extension is
injective: a character that vanishes on a subgroup of finite index
has finite image in the torsion-free group $\Q_\ell$, and so is
zero.  On the other hand, transitivity and the
corestriction--restriction identity give
\[
 \cor_{E/\Q}y
 =\cor_{M/\Q}\cor_{E/M}\res_{E/M}x
 =[E:M]\cor_{M/\Q}x\ne0,
\]
a contradiction.
\end{proof}

\subsection{A signed representation in the character space}

The class field description also controls the supply of characters
when the number field varies.  The regular representation gives upper
bounds for dimensions of subgroup invariants, while the signed
representation below gives lower bounds and ensures a faithful action
on the character space.

\begin{proposition}\label{prop:signed-representation}
Let $M/\Q$ be a finite, totally imaginary Galois extension, put
$R=\Gal(M/\Q)$, and let $c\in R$ be a complex conjugation.
For every rational prime $\ell$, there are $R$-equivariant embeddings
\begin{equation}\label{eq:signed-representation}
 W_M:=\Ind_{\langle c\rangle}^{R}(\mathrm{sign})
 \lhook\joinrel\longrightarrow\X{\ell}{M}
 \lhook\joinrel\longrightarrow\Q_\ell[R],
\end{equation}
where $\mathrm{sign}(c)=-1$.  In particular, the action of $R$ on
$\X{\ell}{M}$ is faithful.
\end{proposition}

\begin{proof}
By the normal basis theorem, $V_M=M\otimes_\Q\Q_\ell$ is the
regular representation of $R$, as is its dual.  The naturality of
\eqref{eq:character-class-field} gives the second embedding.

The equivariant Dirichlet unit theorem
\cite[Section~3]{bass1966} identifies the representation
$\mathcal O_M^\times\otimes_\Z\Q_\ell$ with the permutation
representation on the archimedean places minus its trivial summand.
Since these places are indexed by $R/\langle c\rangle$, writing
$U^{\mathrm{alg}}=\mathcal O_M^\times\otimes_\Z\Q_\ell$ gives
\[
 \Q_\ell\oplus U^{\mathrm{alg}}
 \simeq\Ind_{\langle c\rangle}^{R}\mathbf1.
\]
The character identity follows first over the real numbers from the
usual logarithmic unit embedding, and hence over $\Q_\ell$ by
semisimplicity in characteristic zero.  In particular,
$U^{\mathrm{alg}}$ is self-dual.  Its $\ell$-adic logarithmic image
$U_M$ is a quotient of $U^{\mathrm{alg}}$; no injectivity of the
logarithm on this tensor product is needed.

The decomposition of the regular representation induced from the
two characters of $\langle c\rangle$ is
\[
 \Q_\ell[R]
 \simeq\Ind_{\langle c\rangle}^{R}\mathbf1\oplus W_M
 \simeq\Q_\ell\oplus U^{\mathrm{alg}}\oplus W_M.
\]
Since $U_M^\vee$ is a subrepresentation of
$(U^{\mathrm{alg}})^\vee\simeq U^{\mathrm{alg}}$, semisimplicity
and $\X{\ell}{M}\simeq(V_M/U_M)^\vee$ show that
$\X{\ell}{M}$ contains $\Q_\ell\oplus W_M$.  This proves the
first embedding in \eqref{eq:signed-representation}.

Finally, realize $W_M$ as the direct sum of the signed lines indexed
by the left cosets of $\langle c\rangle$.  An element acting
trivially must fix the line of the identity coset, and therefore
belongs to $\langle c\rangle$.  The nonidentity element $c$ acts
by $-1$ on that line.  Thus only the identity acts trivially on
$W_M$, proving faithfulness.
\end{proof}
